% Procedencia HMT: output/EXTREMA_MEDIA_RAZON_NARRACION_Y_REVISION_20260916/
% ARTICULO/ES/source/libro/01_acoplamiento.tex, lib01:cadena-registro y lib01:K.
% Dirección: ARTICULO/ES/source/sections/k_direccion_dimensional.tex.
% Estatuto de las redes, cartas y flujos de este apéndice: FORMALIZACION_NUEVA.
% La extracción de K y el proyector P_3 son RESULTADO_RECUPERADO.
\section{Medición planar de frontera, coordenadas de Plücker y mutaciones positivas}
\label{app:network-cluster}

La publicación dodecafásica del estado enriquecido admite una realización
positiva cuya red, coordenadas y cambios de carta pueden escribirse
completamente. La composición de partida es
\[
 \mathrm{APP}\longrightarrow\mathrm{TRIT}\longrightarrow\mathrm{TPK}
 \longrightarrow x_{\rm term}\longrightarrow\mathcal R_{12}
 \longrightarrow(b^{90},b^{120},Q_{\rm TPK})
 \longrightarrow U_{\rm sgn}\longrightarrow K.
\]
Esta lectura actúa sobre el estado enriquecido que conserva las dos hojas,
la orientación, el acarreo, la memoria y la frontera. El registro ya generado
es
\[
 K=(234,543,140,729,659,824,621,58,914,794,146,601),
 \qquad Q=\sum_{r=1}^{12}K_r=6263.
\]
Su extracción se recupera exactamente de los dos canales y la carga total:
\[
 b^{90}=(-495,-116,-684,108,601,-90,-173,-88,313,560,-397,461),
\]
\[
 b^{120}=(-425,-281,-481,671,-255,30,475,-543,680,251,6,-128).
\]
En las órbitas de posiciones \((r,r+3,r+6,r+9)\), para \(r=1,2,3\),
sean \(B_r=\sum_{j=0}^3b^{120}_{r+3j}\) y
\(d_{r,j}=b^{90}_{r+3j}\). Se obtienen
\[
 s_1=(Q+2B_1+B_2)/3,\quad s_2=s_1-B_1,\quad s_3=s_2-B_2,
\]
\[
 U^{(r)}=(s_r,d_{r,1}-d_{r,3},d_{r,0}+d_{r,2},d_{r,0}-d_{r,2}),
 \qquad K|_r=\tfrac14H_4U^{(r)},
\]
\[
 H_4=\begin{pmatrix}1&1&1&1\\1&1&-1&-1\\1&-1&1&-1\\1&-1&-1&1\end{pmatrix},
 \qquad H_4^2=4I_4.
\]
La red construida a continuación realiza los valores positivos producidos
por esa extracción. Su forma combinatoria es una formalización posterior
del registro, con origen y orientación declarados.

\subsection{Red dirigida acíclica y determinantes de caminos no intersectantes}

Para un registro positivo de longitud \(L\), ponemos
\[
 d_r=K_r/Q>0,\qquad t_0=0,\qquad
 t_j=\sum_{r=1}^jd_r\quad(1\le j\le L).
\]
Así \(t_L=1\). Las \(L\) separaciones producen \(L+1\) puntos marcados;
en el caso dodecafásico son trece puntos. La matriz que los realiza es
\begin{equation}
 C(d)=\begin{pmatrix}1&1&\cdots&1\\t_0&t_1&\cdots&t_L\end{pmatrix}.
 \label{eq:nc-path-matrix}
\end{equation}
El Grassmanniano \(\operatorname{Gr}(2,L+1)\) identifica matrices de rango
dos por cambios invertibles de las dos filas. Su parte positiva está formada
por las clases que admiten todos sus menores ordenados de orden dos
estrictamente positivos.

Construimos una red dirigida con vértices \(a_j,b_j\), para
\(0\le j\le L\), y sumideros \(q_j\). Hay aristas horizontales
\(a_{j-1}\to a_j\) y \(b_{j-1}\to b_j\), de peso uno; aristas verticales
\(a_j\to b_j\), de peso \(d_j\), para \(j\ge1\); y aristas
\(b_j\to q_j\), de peso uno. Las fuentes ordenadas son \(b_0,a_0\).
Las dos filas se dibujan horizontalmente, las aristas verticales descienden
y los sumideros salen por debajo de la fila inferior. Este dibujo es planar
y la orientación es acíclica. La medición de frontera \(M_{ij}\) es la
suma de los productos de pesos de todos los caminos desde la fuente \(i\)
al sumidero \(q_j\).
La red de caminos tiene dos fuentes y \(L+1\) sumideros, es decir,
\(L+3\) conexiones exteriores. La matriz que realiza tiene \(L+1\)
columnas. La semilla plábica introducida después tiene precisamente
\(L+1\) vértices de frontera y constituye otra presentación de sus
coordenadas; estas dos cuentas de frontera se mantienen explícitas.

\begin{theorem}[Medición de frontera y menores del registro]
\label{thm:nc-network}
La medición de frontera de esta red es \(M=C(d)\). Para
\(0\le i<j\le L\),
\begin{equation}
 \Delta_{ij}(C(d))=t_j-t_i=\sum_{r=i+1}^{j}d_r>0.
 \label{eq:nc-minor-gap}
\end{equation}
El mismo número es la suma de pesos de las parejas de caminos
vértice-disjuntos \(b_0\to q_i\), \(a_0\to q_j\).
\end{theorem}
\begin{proof}
La fuente inferior tiene un único camino a cada sumidero, de peso uno.
Un camino desde \(a_0\) a \(q_j\) desciende exactamente una vez, por una
arista de índice \(1\le r\le j\); su peso es \(d_r\). Esto prueba la
fórmula de la matriz. El camino inferior hasta \(q_i\) ocupa todos los
vértices \(b_0,\ldots,b_i\). Por tanto, un camino de la otra fuente a
\(q_j\) es disjunto de él precisamente cuando desciende en
\(i<r\le j\). Su suma de pesos es el miembro derecho de
\eqref{eq:nc-minor-gap}. La pareja con destinos intercambiados siempre
comparte un vértice. El desarrollo del determinante cancela todas las
parejas que se cruzan y conserva exactamente las disjuntas. Aquí se ha
verificado directamente, para la red concreta, el mecanismo determinantal
de caminos no intersectantes.
\end{proof}

La identificación de las mediciones de redes planas con las celdas del
Grassmanniano no negativo proporciona su reconocimiento geométrico posterior;
véase A. Postnikov,
\href{https://arxiv.org/abs/math/0609764}{\emph{Total positivity,
Grassmannians, and networks}}. La prueba anterior da, además, una fórmula
finita para cada menor sin utilizar un valor metrológico.

La familia \(d_r>0\), \(\sum_rd_r=1\), tiene dimensión \(L-1\).
Su aplicación al Grassmanniano es inyectiva: si las filas de \(C(d)\) y
\(C(d')\) generan el mismo plano, las segundas filas difieren por una
transformación afín; las condiciones \(t_0=t'_0=0\) y
\(t_L=t'_L=1\) fijan esa transformación. En particular, para \(L=12\)
se obtiene una familia de dimensión once dentro de la celda positiva de
dimensión veintidós. El punto está en la celda superior; la familia conserva
su dimensión propia. Los trece puntos y las doce separaciones tienen así
papeles matemáticos distintos.

Permitir \(d_r\ge0\) conserva la no negatividad. Cuando \(d_r=0\), las
columnas consecutivas \(r-1\) y \(r\) coinciden y forman elementos
paralelos del matroide. El rango sigue siendo dos porque \(t_0=0\) y
\(t_L=1\). Las bases son exactamente los pares \(i<j\) que atraviesan
alguna separación positiva. Esta degeneración produce estratos de incidencia
consecutiva. Una columna nula, que daría un lazo, corresponde a una operación
distinta de la coincidencia de dos columnas no nulas.

\subsection{Semilla de triangulación y relación de intercambio}

Sea \(N=L+1\) y considérense los vértices \(0,\ldots,L\) de un polígono
convexo en su orden cíclico. A cada lado o diagonal \((i,j)\) se asigna
la coordenada \(A_{ij}=\Delta_{ij}(C(d))\), con índices ordenados.
La triangulación en abanico desde el vértice cero tiene diagonales
\((0,j)\), \(2\le j\le L-1\). Sus lados son variables congeladas y
sus diagonales variables mutables. El quiver se obtiene orientando
cíclicamente los tres lados de cada triángulo y cancelando los pares de
flechas opuestas. Esta es una semilla de tipo \(A_{N-3}\); para el registro
dodecafásico, su tipo es \(A_{10}\).

La semilla tiene una realización plábica explícita. Se coloca un vértice
negro dentro de cada triángulo; en los vértices del polígono se coloca un
vértice blanco si incide en alguna diagonal, y negro en caso contrario.
Se unen los centros a las tres esquinas de sus triángulos, se retiran los
lados originales y se contraen las aristas entre vértices del mismo color.
Se añade una hoja de frontera a cada vértice marcado. Esta construcción
recoge el algoritmo de triangulación de Kodama y Williams,
\href{https://people.math.harvard.edu/~williams/papers/KW-ArxivVersion.pdf}
{\emph{KP solitons and total positivity for the Grassmannian}, algoritmo 12.1}.
Las etiquetas de coordenadas de esta semilla se evalúan en los menores ya
producidos por la red de dos filas. Para una cara mutable, la coordenada
de tipo \(\mathcal X\) es el cociente de los productos de las
coordenadas \(\mathcal A\) incidentes que prescribe su columna del
quiver; en un cuadrilátero se escribirá explícitamente como una razón
doble. Cambiar el grafo cambia la carta de esos menores, manteniendo su
procedencia.

\begin{theorem}[Intercambio positivo de la realización del registro]
\label{thm:nc-ptolemy}
Para cuatro vértices \(a<b<c<d\),
\begin{equation}
 A_{ac}A_{bd}=A_{ab}A_{cd}+A_{ad}A_{bc}.
 \label{eq:nc-ptolemy}
\end{equation}
En consecuencia, el cambio de diagonal \((a,c)\leftrightarrow(b,d)\)
se evalúa mediante
\[
 A_{bd}=\frac{A_{ab}A_{cd}+A_{ad}A_{bc}}{A_{ac}}>0.
\]
Todas las triangulaciones proporcionan cartas positivas compatibles sobre
la misma realización \(C(d)\).
\end{theorem}
\begin{proof}
Al sustituir \(A_{ij}=t_j-t_i\), la igualdad se reduce a
\[
 (t_c-t_a)(t_d-t_b)
 =(t_b-t_a)(t_d-t_c)+(t_d-t_a)(t_c-t_b).
\]
La expansión de ambos miembros prueba la identidad. Las diferencias son
positivas; por tanto, cada flip conserva una evaluación positiva y el
cociente está definido. La conexión de las triangulaciones por flips
permite pasar entre ellas mediante estos intercambios. La interpretación
cluster de la relación de Plücker es el reconocimiento posterior desarrollado
por J. Scott en
\href{https://arxiv.org/abs/math/0311148}{\emph{Grassmannians and Cluster
Algebras}}. La prueba algebraica de la evaluación concreta ya está contenida
en la identidad mostrada.
\end{proof}

\subsection{Refinamiento de las separaciones y naturalidad de los intercambios}

Sea \(B\ge2\). En profundidad \(n\), cada separación \(d_r\) se divide
en \(B^n\) separaciones consecutivas iguales:
\[
 d^{(n)}_{(r-1)B^n+c}=d_r/B^n,
 \qquad 1\le c\le B^n.
\]
Esto define \(LB^n+1\) puntos y una matriz \(C_n\). Si \(\iota_n(j)=Bj\)
identifica un punto con su descendiente de igual posición, entonces
\begin{equation}
 C_{n+1}|_{\iota_n(\{0,\ldots,LB^n\})}=C_n,
 \qquad
 \Delta_{\iota_n(i),\iota_n(j)}(C_{n+1})=\Delta_{ij}(C_n).
 \label{eq:nc-refinement}
\end{equation}
La red refinada reemplaza cada arista vertical de peso \(d\) por una
secuencia de \(B\) posiciones de descenso, con pesos \(d/B\). La suma
de caminos entre extremos heredados sigue siendo \(d\). La composición de
refinamientos conserva esta igualdad porque las sumas finitas se agrupan
asociativamente.

\begin{proposition}[Naturalidad sobre extremos heredados]
\label{prop:nc-inherited}
Todo intercambio \eqref{eq:nc-ptolemy} entre cuatro extremos heredados
conmuta exactamente con el refinamiento. Toda triangulación antigua se
extiende a una del polígono refinado preservando sus triángulos interiores:
los polígonos añadidos entre dos extremos consecutivos se triangulan por
fuera de esos triángulos. Los flips entre diagonales antiguas conservan
entonces sus relaciones de intercambio.
\end{proposition}
\begin{proof}
La primera conclusión resulta al sustituir cada menor por su igualdad en
\eqref{eq:nc-refinement}. Para la segunda, se conservan las cuerdas que
unen extremos antiguos consecutivos y se triangula cada bolsillo exterior
formado por las nuevas marcas. Estas regiones tienen interiores disjuntos.
Los dos triángulos incidentes en una diagonal antigua permanecen iguales,
y, por tanto, también su cuadrilátero y su identidad de Ptolomeo.
\end{proof}

Esta naturalidad especifica los extremos sobre los que actúa. Un lado
antiguo, al recibir vértices nuevos en su arco de frontera, pasa a ser una
diagonal del polígono mayor. Por ello una variable antiguamente congelada
puede ser mutable en el problema ampliado. La compatibilidad anterior es
una compatibilidad exacta del sistema de evaluaciones y de los flips
heredados. Una inclusión de patrones de semillas con igual régimen de
coeficientes se obtiene manteniendo congeladas esas cuerdas y las nuevas
diagonales auxiliares. Esto declara la operación que conserva la semilla;
el incremento de profundidad y una mutación siguen siendo operaciones
distintas.

\subsection{Acción diagonal proyectiva e invariancia de las razones cruzadas}

Si \(\lambda_j>0\), la acción por columnas
\(C\mapsto C\operatorname{diag}(\lambda_0,\ldots,\lambda_L)\) da
\[
 A_{ij}\longmapsto\lambda_i\lambda_j A_{ij}.
\]
Preserva la positividad y todos los estratos definidos por anulación de
menores. La razón doble asociada a un cuadrilátero es
\begin{equation}
 X_{abcd}=\frac{A_{ab}A_{cd}}{A_{ad}A_{bc}}>0.
 \label{eq:nc-crossratio}
\end{equation}
Los cuatro factores \(\lambda_a\lambda_b\lambda_c\lambda_d\) se
cancelan exactamente. En consecuencia, toda acción diagonal positiva por
columnas deja fijo \(X_{abcd}\), incluso si depende de un parámetro y
tiene determinante uno. La modificación del peso de las columnas conserva
el punto proyectivo de cada columna; sus razones dobles expresan esa
conservación. Para que el registro cambie la forma proyectiva debe actuar
sobre las separaciones ordenadas.

\subsection{Deformación positiva de las separaciones dirigida por el registro}

El corpus fija la representación de las doce posiciones y su proyector
ortogonal \(P_3\) sobre el sector tridimensional. Con
\(P_{11}=I_{12}-J_{12}/12\), la dirección recuperada es
\[
 u=P_3P_{11}K=P_3K\in\mathbf1^\perp.
\]
En la carta declarada, poniendo \(r=\sqrt5\), su evaluación exacta es
\begin{align*}
u=(&-495/4-233r/10,\ 403/4+169r/10,\ -403/4-169r/10,\\
   &495/4+233r/10,\ 601/4+1031r/10,\ 203/4+211r/5,\\
   &-203/4-211r/5,\ -601/4-1031r/10,\ 313/4+569r/10,\\
   &162+219r/20,\ -162-219r/20,\ -313/4-569r/10).
\end{align*}
Se conservan \(\sum_ru_r=0\) y
\(\|u\|^2=(6638585+2275584\sqrt5)/20>0\).

Para \(s\in\mathbb R\), definimos la deformación de los pesos mediante
\begin{equation}
 d_r(s)=\frac{K_re^{su_r}}{\sum_{\ell=1}^{12}K_\ell e^{su_\ell}},
 \qquad t_j(s)=\sum_{r=1}^jd_r(s),\qquad C(s)=C(d(s)).
 \label{eq:nc-gap-flow}
\end{equation}
La fórmula utiliza el registro y su dirección ya generados. Es una nueva
realización analítica del lector; el parámetro \(s\) mide su deformación.
Cuando se usa la salida interna \(\rho_A=\pi_{\rm HMT}/\varphi_{\rm HMT}^2\),
la escritura \(s=\tau\log\rho_A\) expresa exactamente la misma familia.
Esa reparametrización actúa después de la generación de las constantes.
Una identificación de \(s\) o \(\tau\) con un tiempo físico requiere la
ley temporal del lector correspondiente.

\begin{theorem}[Cambio de forma, positividad y compatibilidad con la profundidad]
\label{thm:nc-shape-flow}
La familia \eqref{eq:nc-gap-flow} pertenece a
\(\operatorname{Gr}_{>0}(2,13)\) para todo \(s\) real y conserva
\(t_0=0\), \(t_{12}=1\). Para \(i<j\), sean
\[
 \mu_{ij}(s)=
 \frac{\sum_{r=i+1}^jK_ru_re^{su_r}}{\sum_{r=i+1}^jK_re^{su_r}},
 \qquad
 \mu(s)=\frac{\sum_rK_ru_re^{su_r}}{\sum_rK_re^{su_r}}.
\]
Entonces
\begin{equation}
 \frac{d}{ds}\log A_{ij}(s)=\mu_{ij}(s)-\mu(s),
 \quad
 \frac{d}{ds}\log X_{abcd}(s)
 =\mu_{ab}(s)+\mu_{cd}(s)-\mu_{ad}(s)-\mu_{bc}(s).
 \label{eq:nc-crossratio-derivative}
\end{equation}
La deformación conmuta con cada refinamiento que asigna a los hijos de la
separación \(r\) el peso \(d_r/B\) y la misma velocidad \(u_r\).
\end{theorem}
\begin{proof}
Todos los exponenciales y todos los \(K_r\) son positivos. La normalización
hace que la suma de separaciones sea uno; los menores son sumas de
separaciones positivas por \eqref{eq:nc-minor-gap}. Derivar el logaritmo de
\[
 A_{ij}(s)=\frac{\sum_{r=i+1}^jK_re^{su_r}}
                      {\sum_rK_re^{su_r}}
\]
da la primera identidad. La combinación de cuatro logaritmos en
\eqref{eq:nc-crossratio} cancela las cuatro medias globales con sus signos
y da la segunda. Finalmente, reemplazar \(K_r\) por \(B\) copias de
\(K_r/B\) con idéntica \(u_r\) conserva el numerador y el denominador
cuando se agrupan los hijos. La igualdad vale para todos los parámetros,
para cualquier profundidad y para toda composición de refinamientos.
\end{proof}

Hay un testigo exacto de que esta familia cambia la forma. Para los primeros
cuatro extremos,
\begin{equation}
 X_{0123}(0)=\frac{1560}{23711},\qquad
 \left.\frac{d}{ds}\log X_{0123}(s)\right|_{s=0}
 =-\frac{309899}{917}-\frac{268596}{4585}\sqrt5<0.
 \label{eq:nc-nontrivial-witness}
\end{equation}
La cuenta utiliza únicamente las primeras tres separaciones
\((234,543,140)\) y sus velocidades ya recuperadas. La desigualdad
estricta prueba un movimiento proyectivo efectivo. Las identidades de
Ptolomeo siguen siendo exactas durante ese movimiento, pues se aplican a
las columnas \((1,t_j(s))\) en cada parámetro. El registro controla así
una deformación positiva de la forma y, simultáneamente, un sistema de
cartas que permanece compatible con ella.

El alcance demostrado reúne una red planar de rango dos, su medición
explícita, una semilla plábica de triangulación, todas sus mutaciones por
flips, la naturalidad de los extremos heredados y una deformación de
razones dobles dirigida por \(K\). Cada operación tiene un dominio y un
resultado verificables. La comparación de esta realización con la
dinámica temporal enriquecida y con la rama material utiliza sus respectivos
mapas de lectura; ninguna identidad de coordenadas sustituye esos mapas.
